% ----------------------
\subsection{Section 13 (1838/15 May 1829)}
% ----------------------
	\label{DC13}
	
	% -----
	\subsubsection{Historical Background}
	% -----
		\label{DC13-HB}
		
		Section 13 is an extract from JS's personal history, and was composed in 1838. There is no record contemporaneous to this event which states what happened. 
		
		This description of the priesthood restoration first appeared in the DC in 1876 (HDDC 235).
		
		Here is a more complete extract from \href{https://www.josephsmithpapers.org/paper-summary/history-1838-1856-volume-a-1-23-december-1805-30-august-1834/23}{volume 1-A} of the history:
		
			\begin{quote}
				We still continued the <work of> translation, when in the ensuing month (May, Eighteen hundred and twenty nine) we on a certain day went into the woods to pray and inquire of the Lord respecting baptism for the remission of sins as we found mentioned in the translation of the plates. While we were thus employed praying and calling upon the Lord, a Messenger from heaven, descended in a cloud of light, and having laid his hands upon us, he ordained us, saying unto us; ``Upon you my fellow servants in the name of Messiah I confer the priesthood of Aaron, which holds the keys of the ministring of angels and of the gospel of repentance, and of baptism by immersion for the remission of sins, and this shall never be taken again from the earth, untill the sons of Levi do offer again an offering unto the Lord in righteousness.'' He said this Aaronic priesthood had not the power of laying on of hands, for the gift of the Holy Ghost, but that this should be conferred on <us> hereafter and he commanded us to go and be baptized, and gave us directions that I should baptize Oliver Cowdery, and afterward that he should baptize me. \sout{and that I should be called the first elder of the Church and he the second.} Accordingly we went and were baptized, I baptized him first, and afterwards he baptized me, after which I laid my hands upon his head and ordained him to the Aaronick priesthood, and afterward he laid his hands on me and ordained me to the same priesthood, for so we were commanded.
				
				The messenger who visited us on this occasion and conferred this priesthood upon us said that his name was John, the same that is called John the Baptist in the new Testament, and that he acted under the direction <of> Peter, James, and John, who held the keys of the priesthood of Melchisedeck, which priesthood he said should in due time be conferred on us. And that I should be called the first Elder of the Church and he the second. It was on the fifteenth day of May, Eighteen hundred and twenty-nine that we were baptized; \sout{under} and ordained under the hand of \sout{that} <the> Messenger.
				
				Immediately upon our coming up out of the water after we had been baptized we experienced great and glorious blessings from our Heavenly Father. No sooner had I baptized Oliver Cowdery than the Holy Ghost fell upon him and he stood up and prophecied many things which should shortly come to pass: And again so soon as I had been baptized by him, I also had the Spirit of Prophecy, when standing up I prophecied concerning the rise of this church, and many other things connected with the Church and this generation of the children of men. We were filled with the Holy Ghost, and rejoiced in the God of our Salvation.
				
				Our minds being now enlightened, we began to have the Scriptures laid open to our understandings, and the true meaning and intention of their more mysterious passages revealed unto us, in a manner which we never could attain to previously, nor ever before <had> thought of. In the meantime we were forced to keep secret the circumstances of our having been baptized, and having received this priesthood; owing to a spirit of persecution which had already manifested itself in the neighborhood. We had been threatened with being mobbed, from time to time, and this too by professors of religion, and their intentions of mobbing us, were only counteracted by the influence of my wife's father's family (under Divine Providence) who had became very friendly to me and were opposed to mobs; and were willing that I should be allowed to continue the work of translation without interruption: And therefore offered and promised us protection from all unlawful proceedings as far as in them lay.
			\end{quote}
		
		% --
		\paragraph{HDDC}
		% --
		
			``Joseph Smith and Oliver Cowdery received the Aaronic Priesthood from John the Baptist prior to the events that led to this revelation [Section 11] and also to Section 12. Therefore, these two sections should follow Section 13, which records their ordination'' (note 1, 218).
			
			``According to the Prophet Joseph Smith, his ordination under the hands of John the Baptist preceded the revelation given Hyrum Smith, which revelation is now \hyperref[DC11]{section 11} of the D\&C. Therefore, this section should be placed between sections \hyperref[DC10]{10} and \hyperref[DC11]{11} in order for it to be in its proper chronological order.
			
			``Over ten years after John the Baptist performed the ordinance contained in this revelation, Joseph Smith caused a record of it to be written in his history with this introduction:
			
				\begin{quote}
					``We still continued the work of translation, when in the ensuing month (May, 1829) we on a certain day went into the woods to pray and inquire of the Lord respecting baptism for the remission of sins, that we found mentioned in the translation of the plates. While we were thus employed, praying and calling upon the Lord, a messenger from heaven descended in a cloud of light, and having laid his hands upon us, he ordained us, saying: [Section 13 follows.]
				\end{quote}
				
			\noindent ``However, Oliver Cowdery, who was with the Prophet on this occasion, wrote his account of this ordination five years earlier. This was the first time he had seen an angel, and his statement is filled with awe and wonder. His account is found, as follows, in the \uline{Messenger and Advocate}:
			
				\begin{quote}
					``On a sudden, as from the midst of eternity, the voice of the Redeemer spake peace to us, while the vail was parted and the angel of God came down clothed with glory, and delivered the anxiously looked for message, and the keys of the gospel of repentance!--- What joy! what wonder! what amazement! While the world were racked and distracted---while millions were grouping [sic] as the blind for the fall, and while all men were resting upon uncertainty, as a general mass, our eyes beheld---our ears heard. As in the `blaze of day;' yes, more---above the glitter of the May Sun beam, which then shed its brilliancy over the face of nature! Then his voice, though mild, pierced to the center, and his words, `I am thy fellow servant,' dispelled every fear. We listened---we gazed---we admired! `Twas the voice of the angel from glory---`twas a message from the Most High! and as we heard we rejoiced, while his love enkindled upon our souls, and we were rapt in ths vision of the Almighty! Where was room for doubt? No where: uncertainty had fled, doubt had sunk, no more to rise, while fiction and deception had fled forever!
					
					``But, dear brother think, further think for a moment, what joy filled our hearts and with what surprise we must have bowed, (for who would not have bowed the knee for such a blessing?) when we received under his hand the holy priesthood, as he said, `upon you my fellow servants, in the name of Messiah I confer this priesthood and this authority, which shall remain upon earth, that the sons of Levi may yet offer an offering unto the Lord in righteousness!'%
						\footnote{%
							% \label{}%
							HDDC note: ``\uline{Messenger and Advocate} [Kirtland, Ohio], October, 1834, pp. 15, 16.''
							}
				\end{quote}
			
			``The Prophet's mother also told how this revelation was given, but her account does not agree with that of these two men who were participants. The main thrust of what she wrote is as follows:
			
				\begin{quote}
					``One morning they sat down to work, as usual, and the first thing which presented itself through the Urim and Thummim was a commandment for Joseph and Oliver to repair to the water and attend to the ordinance of baptism. They did so, and as they were returning to the house...%
						\footnote{%
							% \label{}%
							HDDC note: ``Lucy Mack Smith, \uline{History of Joseph Smith by His Mother} (Salt Lake City: Bookcraft, Inc., 1958), p. 142.''
							}%
					(233--234)
				\end{quote}
			
	% -----
	\subsubsection{Versions}
	% -----
		\label{DC13-V}
		
		See the \href{https://drive.google.com/file/d/1goAvNz8jS4R8slYfMG_db55Ty2z_vmgK/view?usp=sharing}{sections comparison} document.
	
		\begin{compactdesc}
			\item[JSP] ---
			\item[BC] ---
			\item[DC 1835] ---
			\item[TS] 3:19 (1 August 1842), 865--866
			\item[MS] 3:9 (January 1843), 148
			\item[DC 1844] ---
		\end{compactdesc}

	% -----
	\subsubsection{Section 13:1} % *DC13-1 
	% -----
		\label{DC13-1}

		UPON you my fellow servants, in the name of Messiah%
			\footnote{%
				% \label{}%
				When I was a missionary, I was once flipping through President Openshaw's scriptures, and noticed that he had pointed out this word in particular---he was interested in the fact that there's no definite article here with ``Messiah.'' It just says ``of Messiah,'' not ``of the Messiah'' or something like that. He took this as a possible indication that ``Messiah'' is a title for a role, and here John is speaking of the person who holds the role as a the holder of the role rather than as this or that particular individual. See \hyperref[CHRIST-role]{this study}.
				}
		I confer the Priesthood of Aaron, which holds the keys of the ministering of angels, and of the gospel of repentance, and of baptism by immersion for the remission of sins;%
			\footnote{%
				% \label{}%
				Compare \hyperref[DC84-26]{DC 84:26--27} and \hyperref[DC107-20]{107:20}.
				}
		and this shall never be taken again from the earth, until%
			\footnote{%
				% \label{}%
				So the Aaronic priesthood \textit{will} be taken again from the earth, after the sons of Levi make their offering. That's the most natural way of reading this passage.
				}
		the sons of Levi%
			\footnote{%
				% \label{}%
				The reference here to the ``sons of Levi'' and an ``offering'' they will make seems to be to \hyperref[MALACHI-3-3]{Malachi 3:3--4}; see notes there. I haven't done an extensive study of this topic, but thinking about the way it's used here in the restoration, it could be referring to the fact that the Lord sees the Levites as no longer able to officiate in their temple duties, which makes it impossible for them to make intervention on behalf of the the people of the house of Israel. For this reason John had to give the priesthood to JS and Cowdery, so that they could begin the process of initiating these ordinances once more, even though they wouldn't naturally be the ones to whom this duty belonged. At some point, those who descend from the ones who possessed this duty will be able to make these offerings again, and at that point, the priesthood will be ``taken again from the earth,'' perhaps to be permanently replaced by the authority of the Melchizedek priesthood.
				
				It is significant that John the Baptist is the one who appears and delivers this priesthood to JS and Cowdery. John was the son of Zechariah or Zacharias, a priest who himself officiated in temple ordinances; John's mother was Elizabeth, noted as being ``of the daughters of Aaron'' (\hyperref[LUKE-1-5]{Luke 1:5}).
				}
		do offer again an offering unto the Lord in righteousness.

		% \begin{framed}
		% \scriptsize
			% \begin{compactdesc}
				% \item[JSP] 
				% % \item[BC] 
				% % \item[]
			% \end{compactdesc}
		% \end{framed}